% A.94 reserved by the root.  Do not include before independent audit.
\subsection{The remaining dyadic contact orders and their multipliers}

\begin{statement}{A.94}{Proved}
{Dyadic inertia determines the small odd contact orders}
Let $s$ have a labeled integral realization $E_\lambda$ as
in A.14, and write $h=h(\lambda)$ for the dyadic contact
order in (A.87.1).
For $h=1$ or $h=3$, the actual ordinary dyadic point has
exactly the same valuation pair as $\lambda$.
The corresponding inertial groups have orders eight
and twenty-four, respectively.

For each of $h=1,2,3$, the same fixed global Weil
multiplier as in A.82 satisfies
\[
                    \Delta_\ell\in(\Z_\ell^\times)^2
                                  \quad\text{for every odd }\ell.
\tag{A.94.1}
\]
For $h=4$, the exact valuation-pair conclusion is A.91,
and one has the additional restriction
\[
                              \Delta_7\in(\Z_7^\times)^2.
\tag{A.94.2}
\]
No assertion about $\Delta_2$ is made in either range.
In particular, a nonsquare odd-coefficient multiplier
can remain only at contact order four or when the
dyadic place is potentially multiplicative and locally
cuspidal; A.34 still excludes it whenever any
potentially multiplicative place is ordinary.
\end{statement}
\begin{proof}
We give both the minimal models and the conductor
check required by the local symplectic criterion.
Let $u\in\Z_2^\times$.  For $h=1,3$ the parameters
at the three labels are
$2^h u$, $1-2^h u$, and $u/2^h$.
At the first two labels the original integral
Legendre equations give the same $c_4,\Delta$;
their $c_6$ have opposite signs.
At infinity, actual rational changes of variables
give the equations in the following table:
\[
\begin{array}{c|c|c|c}
 h&\text{label}&\text{integral model or change}
                    &(v_2c_4,v_2c_6,v_2\Delta)\\ \hline
 1&0,1&E_{2u},\ E_{1-2u}&(4,\ge7,6)\\
 1&\infty&X=4x,\ Y=8y:\ Y^2=X(X-4)(X-2u)
                                     &(6,\ge10,12)\\
 3&0,1&E_{8u},\ E_{1-8u}&(4,6,10)\\
 3&\infty&X=16x,\ Y=64y:\ Y^2=X(X-16)(X-2u)
                                     &(6,9,16).
\end{array}
\tag{A.94.3}
\]
Here $v_2(0)=\infty$ in the two rows allowing $c_6=0$.
For example, the infinity rows have
\[
\begin{array}{c|c|c|c}
 h&c_4&c_6&\Delta\\ \hline
 1&64(u^2-2u+4)&512(u+2)(u-1)(u-4)
                              &2^{12}u^2(u-2)^2\\
 3&64(u^2-8u+64)&512(u+8)(u-4)(u-16)
                              &2^{16}u^2(u-8)^2.
\end{array}
\tag{A.94.4}
\]
The first and third rows of (A.94.3) are minimal
because their discriminant valuations are less
than twelve.  The other two rows are minimal as
well.  Indeed, on any integral Weierstrass equation,
either $a_1$ is odd and $c_4$ is odd, or $a_1$ is
even, $4\mid b_2$, $2\mid b_4$, and
$c_4=b_2^2-24b_4$ is divisible by sixteen.
Thus an integral equation cannot have $v_2(c_4)=2$.
A nonminimal change in a row with $v_2(c_4)=6$
would subtract four from that valuation, producing
the impossible value two (or a negative value
under a larger scaling).
This proves the minimality of all four rows.
For $h=1$ their normalized $c_4$ is three modulo four.
Changing a minimal model by an odd unit multiplies
this normalized invariant by an inverse fourth
power, so the stated residue modulo four is unchanged.

\emph{The exact inertia orders.}
Potential good reduction follows from
$v_2(j)=8-2h>0$.
As in A.87, finite inertia acts faithfully on $E[3]$
and embeds into $\mathrm{SL}_2(\F_3)$.
For $h=1$ the exact cubic resolvent of A.39 is
\[
 4t(1-t)=
 \begin{cases}
 8u(1-2u),&t=2u\text{ or }1-2u,\\
 u(2-u),&t=u/2.
 \end{cases}
\tag{A.94.5}
\]
Both are cubes in $\Q_2$: their valuations are
divisible by three and cubing is an automorphism
of $\Z_2^\times$, by simple-root Hensel lifting.
All three roots of the resolvent therefore lie
in the maximal unramified extension, since
$\mu_3$ is unramified.
The actual pair-partition action of A.39 restricts to
\[
 \mathrm{SL}_2(\F_3)\longrightarrow
       \mathrm{PSL}_2(\F_3)=A_4
                  \longrightarrow A_4/V_4\simeq\Z/3.
\tag{A.94.6}
\]
The first map is the action on the four projective
lines and the second is their action on the three
partitions into two unordered pairs.
Its kernel is the inverse image of $V_4$, namely
the quaternion group of order eight.
Thus the inertia order divides eight.
It is not one because these minimal models are bad.
The complete order-two list in
\cite[\S14.2, equation (14.3)]{LegFreitasKraus2022}
is
\[
 (\ge6,6,6),\quad(4,6,12),\quad
               (\ge8,9,12),\quad(6,9,18),
\tag{A.94.7}
\]
and the order-four list in
\cite[Theorem~6]{LegFreitasKraus2022}
is $(5,8,9)$ or $(7,11,15)$.
Neither list contains either contact-one row of
(A.94.3).  The order is consequently eight.

For $h=3$, the good-reduction discriminant scaling
forces the inertia order to be divisible by three,
using minimal discriminant valuation ten or sixteen.
There is no subgroup of order twelve in
$\mathrm{SL}_2(\F_3)$ by (A.87.7).
The order-three rows in Table~5 of
\cite{LegFreitasKraus2022} have discriminant
valuation four or eight.
In the order-six list of its \S14.3,
valuation ten has $v_2(c_4)\ge6$ and $v_2(c_6)=8$,
whereas valuation sixteen does not occur.
These exact lists exclude both orders for
the contact-three rows of (A.94.3).
The remaining order is twenty-four.

The actual dyadic point is ordinary by A.37 and
has potentially good reduction.
A.16 fixes the two valuation parities.
For an odd contact order these parities fix its
label: the three possibilities are respectively
$(1,0),(0,1),(1,1)$ modulo two.
Only contact orders one and three can remain.
Their distinct inertia orders eight and twenty-four
then force the actual order to equal the realized
order.  This proves the asserted exact valuation
pair at $h=1,3$.

\emph{The conductor and symplectic calculation at order eight.}
We verify the precise case of the primary theorem;
the invariant triple alone will not be used as an
unproved conductor assertion.
Consider any order-eight curve with minimal
invariants
\[
       (v_2c_4,v_2c_6,v_2\Delta)=(4,\ge7,6),
                   \qquad\widetilde c_4\equiv3\pmod4.
\tag{A.94.8}
\]
If its conductor is not $2^6$, Proposition~6 and
Table~5 of \cite{LegFreitasKraus2022} apply and
place it in case $D_a$, with conductor $2^5$.
The alternative conductor $2^6$ is impossible
here, for the following explicit reason.
Its quadratic twist by two has the integral
invariant triple
\[
                   (6,\ge10,12),
                 \qquad\widetilde c_4\equiv3\pmod4.
\tag{A.94.9}
\]
It is minimal by the $c_4$ argument above.
Apply \cite[Theorem~9]{LegFreitasKraus2022} to the
curve paired with itself, at coefficient three.
The preliminary conductor assertion there says
that this two-twist, when the original conductor
is $2^6$, has conductor $2^5$ or $2^8$.
Proposition~6 and Table~5 then apply to it.
For the triple in (A.94.9), their only row is
$D_b$, which requires
$\widetilde c_4\equiv1\pmod4$.
This contradicts its explicitly unchanged unit
residue three.  Thus (A.94.8) always has conductor
$2^5$, and is exactly case (a) of Table~3 in
Theorem~9.

Both curves at labels zero or one and contact
order one satisfy (A.94.8).
At label infinity, apply the same quadratic
two-twist to the two compared curves.
For $t=u/2$ this twist is literally
\[
                     E_{u/2}^{(2)}:
                        y^2=x(x-2)(x-u).
\tag{A.94.10}
\]
It has $c_4=16(u^2-2u+4)$, $v_2(c_6)\ge7$,
and $\Delta=2^6u^2(u-2)^2$.
Thus these two twists also satisfy (A.94.8),
with conductor $2^5$ and the same Table~3 case (a).
Simultaneously twisting the compared representations
preserves the given multiplier: the curve
isomorphisms defining the twists preserve their
Weil pairings, and the same scalar quadratic
character is tensored on both sides.

The actual integral intertwiner identifies the
inertial fields, as in A.87.
For an odd coefficient with $(2/\ell)=1$,
Theorem~8 of \cite{LegFreitasKraus2022}
gives square multiplier.
For $(2/\ell)=-1$, Theorem~9 gives the same
conclusion because the two curves have just
been placed in its same case (a).
Their uniqueness clauses apply to every actual
intertwiner, rather than to the existence of
one specially chosen comparison.
This proves (A.94.1) for $h=1$.

For $h=3$, both curves have the same label and
the same minimal discriminant valuation in
(A.94.3), ten or sixteen.
Their common inertia has order twenty-four.
Theorem~7 of \cite{LegFreitasKraus2022}
therefore gives square multiplier at every
odd coefficient.  The case $h=2$ is A.87.
In every case a square residue of a unit at
an odd coefficient lifts to a square unit.
The local geometric representative changes
and the Weil-paired model identifications
have determinant one, as proved in A.82 and
A.87.  Hence these are the fixed global
$\Delta_\ell$, proving (A.94.1) in its stated scope.

\emph{The coefficient-seven restriction at contact four.}
By A.91 the actual point and the realized
parameter have the same contact-four label.
After an equal quadratic twist they are the
good curves $G_u,G_v$ of (A.91.2) at labels
zero and one, and $G_{1/u},G_{1/v}$ at infinity.
At labels zero and one the common twisting
characters are respectively $\chi_{-1}$ and
the trivial character.
At infinity they are $\chi_{-u}$ and
$\chi_{-v}$; A.16's equality of the first
Kummer squareclasses gives $[u]=[v]$, so
these are the same character on all of
$G_{\Q_2}$, not just on inertia.
The explicit curve maps of A.91 therefore
transport the actual comparison to these
two good curves without changing its
multiplier.

For every odd coefficient, reduction of their
finite etale torsion group schemes gives
Weil-pairing-preserving specialization
isomorphisms to the same curve
\[
                    \overline G:\quad
                    Y^2+XY=X^3+X\quad\text{over }\F_2.
\tag{A.94.11}
\]
This is the actual common reduction, since
each of $u,v$ is an odd unit.
There are four rational points on
$\overline G$: one over $X=0$, two over
$X=1$, and its point at infinity.
Its arithmetic Frobenius $M$ on $T_7$ has
\[
          \operatorname{tr}(M)=-1,\qquad
          \det(M)=2,\qquad
          (\operatorname{tr}M)^2-4\det M=-7.
\tag{A.94.12}
\]
It cannot be scalar modulo seven.
Indeed $M=aI+7B$ would make its discriminant
$49((\operatorname{tr}B)^2-4\det B)$,
contrary to (A.94.12).
Thus its reduction is a non-scalar Jordan
matrix with repeated eigenvalue three.
In a basis for its nilpotent part, its
invertible centralizer consists precisely of
\[
                    \begin{pmatrix}a&b\\0&a\end{pmatrix},
                         \qquad a\in\F_7^\times,
                                      \quad b\in\F_7.
\tag{A.94.13}
\]
Every determinant in this group is $a^2$.

Under the two actual specialization
isomorphisms, the given $T_7$ intertwiner
commutes with this same Frobenius.
Its determinant modulo seven is therefore
a square by (A.94.13).
Hensel lifting gives (A.94.2) for the fixed
global multiplier.  This proves only the
coefficient-seven statement in the contact-four
range, as asserted.

Finally $h\ge5$ is the potentially
multiplicative dyadic range.
If the local section there is ordinary,
A.34 makes every coefficient multiplier
square.  The same holds whenever any odd
potentially multiplicative place is ordinary.
Combining these exact alternatives with
(A.94.1)--(A.94.2) proves the last assertion.
\end{proof}

The minimal-model formulas, the exact
pair-partition kernel, the conductor-case
residue check and the Frobenius centralizer
are audited by
\path{code/legendre_dyadic_multiplier_range_certificate.py},
with retained output in
\path{code/legendre_dyadic_multiplier_range_certificate_output.txt}.
The local criteria used here are Theorems
6--9, Proposition~6, Tables~3 and 5,
and \S\S14.2--14.3 of the explicitly
versioned primary source
\cite{LegFreitasKraus2022}.
